\chapter{\ifproject%
\ifenglish Experimentation and Results\else การทดลองและผลลัพธ์\fi
\else%
\ifenglish System Evaluation\else การประเมินระบบ\fi
\fi}

ในบทนี้จะทดสอบเกี่ยวกับการทำงานในฟังก์ชันหลัก ๆ ของระบบ ทั้งในระดับการจำลอง (Simulation) และการพัฒนาเชิงต้นแบบ (Prototype)

\section{วัตถุประสงค์ของการประเมินระบบ}

การประเมินระบบในโครงงานนี้มีวัตถุประสงค์เพื่อจำลองการทำงานและแสดงแนวคิดเชิงปฏิบัติ (Proof of Concept) ของระบบโหวตความปลอดภัยสูงที่ใช้เทคโนโลยี Quantum Key Distribution (QKD) และ Blind Signature

นอกจากนี้ โครงงานยังมีการพัฒนาในระดับ implementation เบื้องต้น เพื่อแสดงให้เห็นว่าระบบสามารถนำไปใช้งานได้จริงในบางส่วน โดยมุ่งเน้นการทดสอบกลไกหลักของระบบ เช่น การเข้ารหัส การแจกจ่ายกุญแจ และกระบวนการโหวต

การประเมินจะครอบคลุมในประเด็นต่อไปนี้

\begin{itemize}
    \item ความเป็นไปได้ของการใช้ QKD ในการแจกจ่ายกุญแจสำหรับระบบโหวต
    \item ความสามารถของ Blind Signature ในการรักษาความเป็นนิรนามของผู้โหวต
    \item ความถูกต้องของกระบวนการโหวตภายใต้การจำลอง
    \item การแสดงให้เห็นถึงความปลอดภัยของระบบในเชิงแนวคิด
\end{itemize}

\section{ขอบเขตการทดสอบ}

การประเมินระบบแบ่งออกเป็น 2 ส่วนหลัก ได้แก่

\subsection{Simulation Level}

\begin{itemize}
    \item จำลองกระบวนการ QKD (BB84 protocol)
    \item จำลองการตรวจจับการดักฟัง (Eavesdropper)
    \item จำลองการเข้ารหัสและถอดรหัสข้อมูล
    \item จำลองการทำ Blind Signature
\end{itemize}

\subsection{Implementation Level (Prototype)}

\begin{itemize}
    \item พัฒนาระบบโหวตต้นแบบ (เช่น web-based หรือ script)
    \item ทดสอบการส่งข้อมูลโหวตจริงผ่านระบบ
    \item ทดสอบการใช้งานของผู้ใช้ในสถานการณ์จำลอง
\end{itemize}

\section{วิธีการทดสอบ}

\subsection{การทดสอบความถูกต้อง (Functional Testing)}

ทำการทดสอบการทำงานของระบบในแต่ละขั้นตอน เช่น การโหวต การเข้ารหัส และการถอดรหัสข้อมูล เพื่อให้แน่ใจว่าระบบทำงานได้ถูกต้องตามที่ออกแบบไว้

\subsection{การทดสอบความปลอดภัย (Security Testing)}

ทดสอบการป้องกันการโจมตี เช่น

\begin{itemize}
    \item การดักฟังข้อมูล (Eavesdropping)
    \item การปลอมแปลงข้อมูล (Tampering)
    \item การโจมตีแบบ Man-in-the-Middle
\end{itemize}

โดยใช้คุณสมบัติของ QKD ที่สามารถตรวจจับการดักฟังได้